\ifdefined\withModuleSIVP
\else
Let \(d\) be a power of two, and set
\[
 R_d:=\Z[y]/(y^d+1),
 \qquad
 K_d:=\Q[y]/(y^d+1).
\]
We identify elements of \(K_d\) with their degree-\(<d\) representatives and
use the coefficient norm.  On \(K_d^2\), we use the corresponding product
norm.  Since \(y^d+1\) is irreducible over \(\Q\), \(K_d\) is a field.
For a free rank-two \(R_d\)-module
\(\mathcal M\subseteq K_d^2\), let \(\mathcal L(\mathcal M)\subseteq
\mathbb R^{2d}\) denote its \emph{coefficient lattice}.

\paragraph{Technical overview.}
We start with a principal cyclotomic exact-CVP instance
\((g,T,\Delta)\).
Here \(g\in R_d\setminus\{0\}\) generates the ideal \((g)\), \(T\in R_d\) is
the target, and \(\Delta\in\mathbb Z_{\geq0}\) is the squared-distance
threshold.  Choose \(Q\) such that
\(Q^2>\max\{\|T\|_2^2,\|g\|_2^2\}\), and set
\[
 \mathcal M_{Q,g,T}:=R_d(Q,T)+R_d(0,g).
\]
The submodule \(R_d(0,g)\) spans the one-dimensional \(K_d\)-subspace
\(V:=K_d(0,g)\).  Hence any two module vectors linearly independent over
\(K_d\) include a vector outside \(V\).  The choice of \(Q\) ensures that a
shortest such vector can be written as
\(a(Q,T)+b(0,g)\) with \(\|a\|_2^2=1\), hence \(a=\pm y^j\).
Minimizing its second coordinate recovers the
distance from \(T\) to \((g)\), giving
\[
 \bigl(\lambda_2^{K_d}(\mathcal M_{Q,g,T})\bigr)^2
 =Q^2+\dist(T,(g))^2.
\]
Thus the squared-distance gap \(\Delta\) versus \(\Delta+4\) transfers
directly to module-SIVP.
\fi

\ifdefined\withModuleSIVP
\else
For a free rank-two \(R_d\)-module \(\mathcal M\subseteq K_d^2\), define
\[
 \lambda_2^{K_d}(\mathcal M)
 :=\min\left\{r\geq0:
 \dim_{K_d}\Span_{K_d}
 \{\mathbf v\in\mathcal M:\|\mathbf v\|_2\leq r\}\geq2
 \right\}.
\]

We now formulate the corresponding exact decision problem.
\begin{definition}[Cyclotomic exact module-SIVP]
\label{def:module-sivp}
Given a power of two \(d\), an \(R_d\)-basis of a free rank-two
module \(\mathcal M\subseteq K_d^2\), and a nonnegative integer squared
threshold \(\Gamma\), exact decision-module-SIVP asks whether
\[
 \bigl(\lambda_2^{K_d}(\mathcal M)\bigr)^2\leq\Gamma.
\]
\end{definition}
\fi

The following lemma gives the embedding used in the reduction
and identifies the second \(K_d\)-linear successive minimum.

\begin{lemma}
\label{lem:module-completion}
Let \(g\in R_d\setminus\{0\}\), let \(T\in R_d\), and choose an integer
\(Q>0\) such that
\[
 Q^2>\max\{\|T\|_2^2,\|g\|_2^2\}.
\]
Set
\[
 \mathcal M_{Q,g,T}
 :=R_d(Q,T)+R_d(0,g)
 =\{(Qa,Ta+gb):a,b\in R_d\}.
\]
Then \(\mathcal M_{Q,g,T}\) is a free rank-two \(R_d\)-module
with {basis}
\[
 \begin{pmatrix}
  Q&0\\
  T&g
 \end{pmatrix},
\]
and
\[
 \bigl(\lambda_2^{K_d}(\mathcal M_{Q,g,T})\bigr)^2
 =Q^2+\dist(T,(g))^2.
\]
\end{lemma}

\begin{proof}
The matrix in the statement has determinant \(Qg\neq0\) in
\(K_d\), so \((Q,T)\) and \((0,g)\) are \(K_d\)-linearly independent.
Since these vectors span \(\mathcal M_{Q,g,T}\) over \(R_d\), they form an
\(R_d\)-basis, and hence \(\mathcal M_{Q,g,T}\) is free of rank two.

Let \(V:=K_d(0,g)\), and set
\[
 r^2:=Q^2+\dist(T,(g))^2.
\]
Every vector of \(\mathcal M_{Q,g,T}\setminus V\) has the form
\(\mathbf v(a,b):=(Qa,Ta+gb)\) with \(a\neq0\).
Since \(a\in R_d\setminus\{0\}\), \(\|a\|_2^2\) is a positive integer.

\textit{Case 1: \(\|a\|_2^2=1\).}
Then \(a=\pm y^j\) for some \(0\leq j<d\).  Multiplication by \(a\) is an
isometry and \(a^{-1}\in R_d\).  Therefore
\[
 \|\mathbf v(a,b)\|_2^2
 =Q^2+\|T+g(ba^{-1})\|_2^2
 \geq r^2.
\]

\textit{Case 2: \(\|a\|_2^2\geq2\).}
In this case,
\[
 \|\mathbf v(a,b)\|_2^2
 \geq2Q^2
 >Q^2+\dist(T,(g))^2
 =r^2,
\]
where \(\dist(T,(g))\leq\|T\|_2<Q\).  Thus every vector outside \(V\) has
norm at least \(r\).

Let \(c\in(g)\) be a closest element to \(T\).  Then
\((Q,T-c)\in\mathcal M_{Q,g,T}\) has norm \(r\) and lies outside \(V\), while
\(\|(0,g)\|_2<Q\leq r\).  Since \((0,g)\) and \((Q,T-c)\) are linearly
independent over \(K_d\),
\[
 \lambda_2^{K_d}(\mathcal M_{Q,g,T})=r.
\]
\end{proof}

We now apply Lemma~\ref{lem:module-completion} to the
principal-ideal CVP instances from Theorem~\ref{thm:cyclotomic-cvp}.

\Needspace{12\baselineskip}
\begin{theorem}
\label{thm:module-sivp}
Exact decision-module-SIVP is \(\NP\)-complete on free rank-two modules over
the power-of-two cyclotomic rings \(R_d\).
Hardness holds even for lower-triangular \(R_d\)-bases of the form
\[
 \begin{pmatrix}
  Q&0\\
  T&g
 \end{pmatrix},
 \qquad
 Q\in\mathbb{N}^*,\quad g,T\in R_d,\quad g\neq0.
\]
\end{theorem}

\begin{proof}
Apply Theorem~\ref{thm:cyclotomic-cvp} to obtain a principal-ideal CVP
instance \((d,g,T,\Delta)\) satisfying
\[
\begin{aligned}
 \dist(T,(g))^2
   &=\Delta &&\text{in a YES instance},\\
 \dist(T,(g))^2
   &\geq\Delta+4 &&\text{in a NO instance}.
\end{aligned}
\]
Set
\[
 Q:=1+\left\lceil
 \sqrt{\max\{\|T\|_2^2,\|g\|_2^2\}}
 \right\rceil,
 \qquad
 \mathcal M:=\mathcal M_{Q,g,T},
 \qquad
 \Gamma:=Q^2+\Delta.
\]
Lemma~\ref{lem:module-completion} gives
\[
\begin{aligned}
 \bigl(\lambda_2^{K_d}(\mathcal M)\bigr)^2
   &=\Gamma &&\text{in a YES instance},\\
 \bigl(\lambda_2^{K_d}(\mathcal M)\bigr)^2
   &\geq\Gamma+4 &&\text{in a NO instance}.
\end{aligned}
\]
The basis of \(\mathcal M\) and the threshold \(\Gamma\) have
polynomially bounded encoding length, so the construction is a polynomial-time
reduction.

For membership in \(\NP\), a YES certificate consists of two vectors in
\(\mathcal M\) of squared norm at most \(\Gamma\).
The verifier checks membership of both coefficient vectors in
\(\mathcal L(\mathcal M)\), their norm bounds, and \(K_d\)-linear independence
by testing whether the corresponding \(2\times2\) determinant is nonzero.
After clearing the denominators of the
input basis, the norm bound gives certificates of polynomial bit length.
\end{proof}

\ifdefined\withModuleSIVP
{By the componentwise form of Lemma~\ref{lem:canonical}, the
theorem also holds under the componentwise canonical embedding after
multiplying squared thresholds by \(d\).}
\fi
